Der Luftdruck auf Meereshöhe, $p_0 = \pLO$ (Normaldruck, FoTaBe S.~186), ist der Schweredruck der ganzen Luftsäule darüber. Hätte die Luft überall die Dichte $\rho_0 = \rhoLO$ (bei Normalbedingungen, FoTaBe S.~187), so wäre $p_0 = \rho_0\,g\,h$ und
\begin{align}
 h &= \hAF \\
   &= \frac{\pL}{\rhoL\cdot\gE} \\
   &= \hA \approx \result{\hAP{3}{2}}
\end{align}
In Wirklichkeit nimmt die Dichte mit der Höhe ab (Luft ist kompressibel); die Atmosphäre reicht deshalb viel höher, wird aber immer dünner.
